\documentclass[11pt,a4paper,UTF8]{ctexart}
\usepackage{geometry}
\geometry{left=18mm,right=18mm,top=24mm,bottom=24mm,headheight=22pt,footskip=12mm}
\usepackage{amsmath,amssymb,mathtools,bm,esint,extarrows,centernot}
\usepackage{xcolor}
\usepackage{tcolorbox}
\tcbuselibrary{skins,breakable}
\usepackage{fancyhdr}
\usepackage{titlesec}
\usepackage{hyperref}
\usepackage{tikz}
\usepackage{booktabs}
\usepackage{tabularx}

% --- Color Palette (Purple & DeepPink Academic Theme) ---
\definecolor{DeepPurple}{HTML}{4C1D95}
\definecolor{TitlePurple}{HTML}{581C87}
\definecolor{SubPurple}{HTML}{6B21A8}
\definecolor{DeepPink}{HTML}{FF1493}
\definecolor{LightPinkBg}{HTML}{FFF5F8}
\definecolor{LightPurpleBg}{HTML}{F8F5FF}
\definecolor{GoldAccent}{HTML}{D97706}
\definecolor{DarkText}{HTML}{1F2937}

\hypersetup{
    colorlinks=true,
    linkcolor=TitlePurple,
    citecolor=DeepPink,
    urlcolor=SubPurple,
    pdfborder={0 0 0}
}

% --- Typography & Spacing ---
\linespread{1.3}
\setlength{\parskip}{0.35em plus 0.1em minus 0.1em}
\setlength{\parindent}{0pt}

% --- Section Titles Styling ---
\titleformat{\section}
  {\Large\bfseries\color{TitlePurple}}{\thesection}{1em}{}
  [\vspace{1mm}{\color{TitlePurple!30}\hrule height 1pt}]
\titleformat{\subsection}
  {\large\bfseries\color{SubPurple}}{\thesubsection}{0.8em}{}
\titleformat{\subsubsection}
  {\normalsize\bfseries\color{DeepPurple}}{\thesubsubsection}{0.6em}{}

% --- tcolorbox Environments ---
% 1. Theorem / Formula / Method / Analysis Box (Deep Purple)
\newtcolorbox{thmbox}[1][]{
    enhanced,
    breakable,
    colback=LightPurpleBg,
    colframe=TitlePurple,
    arc=3pt,
    boxrule=0pt,
    leftrule=3.5pt,
    left=10pt,
    right=10pt,
    top=8pt,
    bottom=8pt,
    title=#1,
    coltitle=white,
    colbacktitle=TitlePurple,
    fonttitle=\bfseries\small,
    attach boxed title to top left={yshift=-2mm, xshift=4mm},
    boxed title style={boxrule=0pt, arc=2pt}
}

% 2. Solution / Formula / Derivation Box (DeepPink)
\newtcolorbox{solbox}[1][]{
    enhanced,
    breakable,
    colback=LightPinkBg,
    colframe=DeepPink,
    arc=3pt,
    boxrule=0pt,
    leftrule=3.5pt,
    left=10pt,
    right=10pt,
    top=8pt,
    bottom=8pt,
    title=#1,
    coltitle=white,
    colbacktitle=DeepPink,
    fonttitle=\bfseries\small,
    attach boxed title to top left={yshift=-2mm, xshift=4mm},
    boxed title style={boxrule=0pt, arc=2pt}
}

% 3. Learning Goals / Summary Box
\newtcolorbox{targetbox}[1][]{
    enhanced,
    breakable,
    colback=LightPurpleBg,
    colframe=DeepPurple,
    arc=3pt,
    boxrule=1pt,
    left=10pt,
    right=10pt,
    top=8pt,
    bottom=8pt,
    title=#1,
    coltitle=white,
    colbacktitle=DeepPurple,
    fonttitle=\bfseries\small,
    attach boxed title to top left={yshift=-2mm, xshift=4mm},
    boxed title style={boxrule=0pt, arc=2pt}
}

\begin{document}

% ------------------- 讲义封面 (Cover Page) -------------------
\begin{titlepage}
\thispagestyle{empty}
\begin{tikzpicture}[remember picture,overlay]
    \fill[DeepPurple] (current page.north west) rectangle ([yshift=-7cm]current page.north east);
    \draw[DeepPink, line width=3.5pt] ([yshift=-7cm]current page.north west) -- ([yshift=-7cm]current page.north east);
    \draw[GoldAccent, line width=1pt] ([yshift=-7.12cm]current page.north west) -- ([yshift=-7.12cm]current page.north east);
    
    \node[anchor=north east, opacity=0.12, text=white] at ([xshift=-1.5cm, yshift=-0.5cm]current page.north east) {
        \fontsize{70}{70}\selectfont\bfseries 2027
    };
\end{tikzpicture}

\vspace*{0.8cm}
\begin{center}
    {\color{white}\fontsize{26}{32}\selectfont\bfseries 2027 考研数学(二) 核心讲义}\\[0.6cm]
    {\color{white}\fontsize{13}{18}\selectfont\bfseries 全国硕士研究生招生考试 · 统考数学(二) 核心精讲全书}\\[1.5cm]
\end{center}

\vspace{3.5cm}
\begin{center}
    {\color{GoldAccent}\large\bfseries $\blacktriangleright$ 基础强化 · 核心题解 · 考点深水区 $\blacktriangleleft$}\\[0.8cm]
    {\color{TitlePurple}\fontsize{24}{28}\selectfont\bfseries 第三章 一元函数微分学的概念}\\[0.6cm]
    {\color{DeepPink}\large\bfseries —— 体系化理论精讲与公式全景 ——}\\[1.5cm]
\end{center}

\vfill

\begin{center}
\begin{tcolorbox}[
    colback=white,
    colframe=DeepPurple,
    width=0.88\textwidth,
    arc=4pt,
    boxrule=1.2pt,
    halign=center,
    drop shadow=gray!30
]
    \vspace{3mm}
    {\bfseries\large 编著团队：EscoffierZhou}\\[2.5mm]
    {\bfseries\color{TitlePurple} 目标院校：中国科学院大学 (UCAS) 计算机专硕 (22408)}\\[2.5mm]
    {\small\color{gray} 备考铁律：手算真题数字 · 错题白纸重做 · 408 客观题为王}\\[2.5mm]
    {\small 编制版本：2027 届首发版 · \today}
    \vspace{2mm}
\end{tcolorbox}
\end{center}
\vspace{0.8cm}
\end{titlepage}

% ------------------- 目录与页面主体配置 -------------------
\pagestyle{fancy}
\fancyhf{}
\fancyhead[L]{\small\color{gray}2027 考研数学(二) · 核心讲义系列}
\fancyhead[R]{\small\color{TitlePurple}\nouppercase{\leftmark}}
\fancyfoot[C]{\small\color{gray}— 第 \thepage\ 页 —}
\fancyfoot[R]{\small\color{gray}UCAS 22408}
\renewcommand{\headrulewidth}{0.5pt}

\pagenumbering{Roman}
\tableofcontents
\newpage
\pagenumbering{arabic}



\section{1.导数的概念与几何意义}

特点:导数是局部概念,研究自变量变化极其微小时因变量的变化速率

\begin{thmbox}[原理[1]: 导数定义(增量式与函数式)]

设$y=f(x)$在$x_0$邻域有定义,极限存在则称在$x_0$处可导:

  增量式:$f'(x_0) = \lim_{\Delta x \to 0} \frac{f(x_0+\Delta x) - f(x_0)}{\Delta x} = \lim_{\Delta \to 0} \frac{f(x_0+\Delta) - f(x_0)}{\Delta}$

  点函数式:$f'(x_0) = \lim_{x \to x_0} \frac{f(x) - f(x_0)}{x - x_0}$

  解释:分母与分子增量必须保持一致且同趋于0,本质是定点与动点的连线斜率在逼近过程中的极限.

可导:	[1]$\lim_{\Delta x \to 0} \frac{\Delta y}{\Delta x}$存在			$\to$[函数增量$\Delta y$和自变量增量$\Delta x$的比值在$\Delta\to0$时候存在]

\textbf{[2]} (称该极限为y=f(x)在$x_0$处的导数,记作
\begin{equation*}
f'(x_0)
\end{equation*}
)

设$u=\lim_{\Delta x \to 0} \frac{\Delta y}{\Delta x}$(比阶),则[u=0,水平切线];[u=k,常数];[u=$\infty$,不可导,但可能存在垂直切线]

\end{thmbox}
\begin{thmbox}[原理[2]: 单侧导数与可导充要条件]

  左导数:$f'_-(x_0) = \lim_{\Delta x \to 0^-} \frac{f(x_0+\Delta x) - f(x_0)}{\Delta x} = \lim_{x \to x_0^-} \frac{f(x) - f(x_0)}{x - x_0}$

  右导数:$f'_+(x_0) = \lim_{\Delta x \to 0^+} \frac{f(x_0+\Delta x) - f(x_0)}{\Delta x} = \lim_{x \to x_0^+} \frac{f(x) - f(x_0)}{x - x_0}$

  充要条件:$f'(x_0) \text{ 存在} \iff f'_-(x_0) \text{ 与 } f'_+(x_0) \text{ 均存在且相等}$.

\end{thmbox}
\begin{thmbox}[原理[3]: 连续与可导的关系(连续;可导;连续可导)]

可导$\Rightarrow$连续; 连续$\not\Rightarrow$可导;不连续$\Rightarrow$不可导

  (反例:$f(x)=|x|$在$x=0$处连续但不可导[尖点])

  判断连续:$\lim\limits_{x \to x_0^+} f(x)=\lim\limits_{x \to x_0^-} f(x)= f(x_0)$				   (极限值=原函数函数值)

  判断可导:
\begin{equation*}
\lim_{x \to x_0^+} \frac{f(x) - f(x_0)}{x - x_0}= \lim_{x \to x_0^-} \frac{f(x) - f(x_0)}{x - x_0}=A(\text{而非}\infty)
\end{equation*}
		 (左差商=右差商)

  连续可导:$\lim_{x \to x_0^+} \frac{f(x) - f(x_0)}{x - x_0}= \lim_{x \to x_0^-} \frac{f(x) - f(x_0)}{x - x_0} = f'(x_0)$		    (先可导,然后判断导函数连续)

\end{thmbox}
\begin{thmbox}[原理[4]: 可导函数的奇偶性/周期性/有界性]

1.奇偶性互换:$f(x)$为可导偶函数     $\Rightarrow$ $f'(x)$为奇函数

	   $f(x)$为可导奇函数      $\Rightarrow$ $f'(x)$为偶函数.

  推论: 若奇函数$g(x)$在$x=0$处可导,必有$g'(x)$为偶函数且$g(0)=0$[每求导一次奇偶性反转一次]

2.周期守恒:$f(x)$为可导且周期为$T$的函数$\Rightarrow$ $f'(x)$亦为周期为$T$的周期函数.

3.有界性\textbf{不传递}:$f(x)$有界$\not\Rightarrow f'(x)$有界

  (反例:$f(x)=\sqrt{x}, x \in (0,1]$有界, 但$f'(x)=\frac{1}{2\sqrt{x}}$无界)

\end{thmbox}
\begin{thmbox}[原理[5]: 导数的几何意义]

函数$y = f(x)$在点$(x_0, y_0)$处的导数值$f'(x_0)$即为曲线在该点处的切线斜率$k$

  切线方程:
\begin{equation*}
y - y_0 = f'(x_0)(x - x_0)
\end{equation*}


  法线方程:
\begin{equation*}
y - y_0 = -\frac{1}{f'(x_0)}(x - x_0)
\end{equation*}
$(f'(x_0) \ne 0)$

\textit{}不是所有函数都有切线(在$x_0$处):尖点/不连续点/剧烈震荡点\textit{}

\textbf{有切线不一定有导数(垂线),有导数必然有切线}

  $y = |x|$在$x=0$处,$f'_+(0) = 1,\ f'_-(0) = -1 \implies$出现尖点,不可导且不存在切线 

  在$x=0$处,$f'(0) = \lim\limits_{x \to 0} \frac{x^{\frac{1}{3}} - 0}{x} = \infty \implies$导数不存在,但切线方程存在（切线为垂直线$x = 0$）。 

\end{thmbox}
\begin{thmbox}[原理[6]: 高阶导数]

[定义]二阶导数:	
\begin{equation*}
f''(x_0) = \lim_{\Delta x \to 0} \frac{f'(x_0 + \Delta x) - f'(x_0)}{\Delta x} = \lim_{x \to x_0} \frac{f'(x) - f'(x_0)}{x - x_0}
\end{equation*}


[定义]$n$阶导数:	
\begin{equation*}
f^{(n)}(x_0) = \lim_{\Delta x \to 0} \frac{f^{(n-1)}(x_0 + \Delta x) - f^{(n-1)}(x_0)}{\Delta x} = \lim_{x \to x_0} \frac{f^{(n-1)}(x) - f^{(n-1)}(x_0)}{x - x_0}
\end{equation*}


\textbf{连续性与阶数的传递性质}

性质(1):[二阶推一阶]若$f(x)$在$x_0$处存在二阶导数$f''(x_0)$，则$f(x)$在$x_0$的某邻域内一阶导数$f'(x)$存在，且$f'(x)$与$f(x)$都在$x_0$处\textbf{连续}

性质(2):[n阶推n-1阶]若$f(x)$在$x_0$处存在$n$阶导数$f^{(n)}(x_0)$，则在$x_0$的某邻域内$1 \sim (n-1)$阶导数均存在，且$0 \sim (n-1)$阶导数在$x_0$处均\textbf{连续}

  即最高阶$f^{(n)}(x)$在邻域内不一定存在，在$x_0$处也不一定连续

\textbf{极值第二充分条件}

设$f(x)$在$x_0$处二阶可导，且$f'(x_0) = 0$（驻点）:

  [$f''(x_0) < 0$]$\implies f(x_0)$为极大值

  [$f''(x_0) > 0$]$\implies f(x_0)$为极小值

  [$f''(x_0) = 0$]$\implies$该判据失效

\end{thmbox}
\begin{thmbox}[原理[7]: 微分的概念和计算]

定义:设函数$y = f(x)$在$x_0$的某邻域内有定义，若自变量增量$\Delta x$对应的函数增量$\Delta y$可表示为：

\end{thmbox}

\begin{equation*}
\Delta y = f(x_0 + \Delta x) - f(x_0) = A \Delta x + o(\Delta x)
\end{equation*}

\begin{solbox}[分步推导与答题规范]
  其中$A$是与$\Delta x$无关的常数,$o(\Delta x)$是比$\Delta x$高阶的无穷小，则称$f(x)$在$x_0$处可微
\end{solbox}

微分记作:$dy = A \Delta x = A dx$(称为函数增量$\Delta y$的线性主部)

对于一元函数：可微$\Leftrightarrow$可导	     且必有$A = f'(x_0)$,故$dy = f'(x) dx$（“以直代曲”）。

对于多元函数：
\begin{equation*}
\text{偏导存在且连续} \implies \text{可微} \implies \text{偏导存在} \quad (\text{双向均不可逆})
\end{equation*}


[可微性判定]写增量:	   $\Delta y = f(x_0 + \Delta x) - f(x_0)$。

[可微性判定]写线性增量:	确定待定系数 $A = f'(x_0)$，线性部分为 $A \Delta x$。

[可微性判定]作差求极限:	
\begin{equation*}
\lim_{\Delta x \to 0} \frac{\Delta y - A \Delta x}{\Delta x} = 0 \iff \Delta y - A \Delta x = o(\Delta x) \iff \text{可微}
\end{equation*}


[几何意义]$\Delta x = dx$（自变量的增量等于自变量的微分）。

[几何意义]$dy = f'(x_0) \Delta x$：切线在纵坐标上的增量（线性主部）。

[几何意义]$\Delta y$：曲线上实际点的真实纵坐标增量。

[几何意义]误差：$\Delta y - dy = o(\Delta x)$（曲线上点与切线点的高度差）。


\section{2.题型}

\begin{thmbox}[题型(1):求高阶导数]

\end{thmbox}

\begin{tcolorbox}[colback=LightPurpleBg,colframe=TitlePurple!60,arc=2pt,boxrule=0.8pt]
求 f\textasciicircum{}(n)(0) (特殊点)\\[1mm]
对称/奇偶? 				[方法A: 奇偶性]\\[1mm]
可化为常见展开? 	 	  [方法B: 泰勒系数](频率最高,算量最小)\\[1mm]
反三角/根式? 	    	[方法E: 微分方程递推]\\[1mm]
求 f\textasciicircum{}(n)(x) (任意点通项)\\[1mm]
乘积且含多项式?	 	[方法C: 莱布尼茨]\\[1mm]
可裂项/降幂		 	  [方法D: 基本公式]\\[1mm]
找规律				 	[猜想+数学归纳法]\\[1mm]
\vspace{1mm}
\end{tcolorbox}

\textit{}(A)求$f^(n)(0)$(特殊点)\textit{}   

\textcolor{DeepPink}{方法A:[在x=0处奇函数(0)=0]根据可导函数的奇偶性直接写答案}

\begin{solbox}[分步推导与答题规范]
操作A:化简,直接看奇偶性,写$f^{(n)}(0)$[偶阶导]答案为0

操作B:构造奇函数拆分(计算$f(x) + f(-x)$)

  理论(1):如果是奇函数,则f(x)恒可唯一分解为一个偶函数与一个奇函数之和

  理论(2):若函数满足$f(x) + f(-x) = 2C$,对称中心$(0, C)$

  操作(1):计算$f(x) + f(-x)$:

  	如果[$f(x) + f(-x) = 0$]$\Rightarrow$锁定对称中心$(0, 0)$		  [f(x)为严格奇函数]

  	如果[$f(x) + f(-x) = 2C$]$\Rightarrow$锁定对称中心$(0, C)$,令$g(x) = f(x) - C$[g(x)为严格奇函数]

  	因为
\begin{equation*}
f^{(n)}(x) = g^{(n)}(x) (n \ge 1)
\end{equation*}
,因此对两端同时求导		[对右侧求导]

  	回到操作A,得到$f^{(n)}(0)$[偶阶导]答案为0
\end{solbox}

\textcolor{DeepPink}{方法B:[在x=0处可以泰勒展开]泰勒展开计算}

\begin{solbox}[分步推导与答题规范]
操作(1):展开目标函数,展开为x的幂级数:	$\sin x, \cos x, e^x, \ln(1+x), \frac{1}{1-x}, (1+x)^\alpha$等

操作(2):定位通项系数:    	          找到$x^n$项的系数$a_n$

操作(3):乘上阶乘:			$f^{(n)}(0) = n! \cdot a_n$  [注:即乘上泰勒公式的分母]

Eg:设$f(x) = \frac{x}{1 + 2x^2}$,求$f^{(9)}(0)$与$f^{(10)}(0)$			   [遇到分式不能直接求导]

本能先看奇偶,是一个奇函数,所以$f^{(10)}(0)=0$

根据泰勒公式:
\begin{equation*}
\frac{1}{1-u} = 1 + u + u^2 + u^3 + u^4 + \cdots + o(u^k)
\end{equation*}


$\Rightarrow\frac{1}{1 - (-2x^2)}$ $= 1 + (-2x^2) + (-2x^2)^2 + (-2x^2)^3 + (-2x^2)^4 + \cdots$ $= 1 - 2x^2 + 4x^4 - 8x^6 + 16x^8 + \cdots$

$\Rightarrow \frac{x}{1 + 2x^2}$ $= x \cdot \left(1 - 2x^2 + 4x^4 - 8x^6 + 16x^8 + \cdots\right)$       $= x - 2x^3 + 4x^5 - 8x^7 + 16x^9 + \cdots$

结论:$f^{(9)}(0)$中$x^9$项的系数为16,所以$f^{(9)}(0)=16\cdot 9!$
\end{solbox}

\textcolor{DeepPink}{方法E:[]微分方程}

\begin{solbox}[分步推导与答题规范]
操作(1):求低阶导数去分母/去根号:	   得到类似$(1-x^2)y' = 1$或$(1-x^2)y'' - x y' = 0$

操作(2):两端应用莱布尼茨公式求n阶导

操作(3):代入x=0求递推关系式

操作(4):整合等式+归纳法奇偶分类

\end{solbox}

Eg:设$y = (\arcsin x)^2$,求$y^{(n)}(0)$

\begin{solbox}[分步推导与答题规范]

\begin{equation*}
y' = 2\arcsin x \cdot \frac{1}{\sqrt{1-x^2}} \implies \sqrt{1-x^2} y' = 2\arcsin x
\end{equation*}


[化简]两边平方去根号:    	
\begin{equation*}
(1-x^2)(y')^2 = 4(\arcsin x)^2 = 4y
\end{equation*}


[化简]两边再次对$x$求导降次:	
\begin{equation*}
-2x(y')^2 + (1-x^2) \cdot 2y'y'' = 4y'
\end{equation*}


[化简]两边约去$2y'$:		
\begin{equation*}
(1-x^2)y'' - xy' = 2
\end{equation*}


[两边同时求导]:对第一项$[(1-x^2)y'']^{(n)}$展开(截断为3项):


\begin{equation*}
[(1-x^2)y'']^{(n)} = (1-x^2)y^{(n+2)} + C_n^1 (-2x) y^{(n+1)} + C_n^2 (-2) y^{(n)}
\end{equation*}
	$= (1-x^2)y^{(n+2)} - 2nx y^{(n+1)} - n(n-1)y^{(n)}$

[两边同时求导]:对第二项$[xy']^{(n)}$展开(截断为2项)


\begin{equation*}
[xy']^{(n)} = x y^{(n+1)} + C_n^1 (1) y^{(n)}
\end{equation*}
				$= x y^{(n+1)} + n y^{(n)}$


\begin{equation*}
(2)^{(n)} = 0
\end{equation*}


[整合等式]:	  
\begin{equation*}
(1-x^2)y^{(n+2)} - (2n+1)x y^{(n+1)} - n^2 y^{(n)} = 0
\end{equation*}


令$x = 0$,中间项含$x$归零:
\begin{equation*}
y^{(n+2)}(0) = n^2 y^{(n)}(0) (n \ge 1)
\end{equation*}


[数学归纳法]:由原式及低阶导数得:$y(0) = 0,\ y'(0) = 0,\ y''(0) = 2$

当$n$为奇数时:		所有奇数阶导数$y^{(2k-1)}(0) = 0$

当$n$为偶数时(令$n = 2k$): 
\begin{equation*}
y^{(2k)}(0) = (2k-2)^2 \cdot y^{(2k-2)}(0)= (2k-2)^2 \cdot (2k-4)^2 \cdots 2^2 \cdot y''(0)
\end{equation*}


			
\begin{equation*}
= [2^{k-1} (k-1)!]^2 \cdot 2 = 2^{2k-1} [(k-1)!]^2
\end{equation*}

\end{solbox}

\textit{}(B)求$f^(n)(x)$(任意点通项)\textit{}

\textcolor{DeepPink}{方法C:[两项乘积,在求导若干次后,其中一项趋于0]莱布尼茨公式}

\begin{solbox}[分步推导与答题规范]
操作(1):令$v(x)$为有限次求导归零的多项式	(如$x^2, x^3$等)

令$u(x)$为易求任意阶导数的函数	 (如$e^{ax}, \sin ax, \cos ax, \frac{1}{ax+b}$)

操作(2):若$v(x)$是$m$次多项式,公式中$k > m$的项全为0,只需计算前$m+1$项

操作(3):若求$f^{(n)}(0)$，在展开式末尾直接代入$x=0$

Eg:设$f(x) = (x^2 - x) e^{-2x}$,求$f^{(n)}(x)$,并计算$f^{(100)}(0)$	[多项式(可截断) 乘以 指数函数]

令$u(x) = e^{-2x}$  $\to u^{(k)}(x) = (-2)^k e^{-2x}$

令$v(x) = x^2 - x$ $\to v'(x) = 2x - 1,\ v''(x) = 2,\ v'''(x) = 0$

根据莱布尼茨公式: 	
\begin{equation*}
f^{(n)}(x) = [C^0_n \cdot (-2)^n e^{-2x} (x^2 - x)] + [C^1_n \cdot (-2)^{n-1} e^{-2x} (2x - 1)] + [C^2_n \cdot (-2)^{n-2} e^{-2x} \cdot 2]
\end{equation*}


		   $= [1 \cdot (-2)^n e^{-2x} (x^2 - x)] + [n \cdot (-2)^{n-1} e^{-2x} (2x - 1)] + [\frac{n(n-1)}{1\cdot 2} \cdot (-2)^{n-2} e^{-2x} \cdot 2]$

	      	
\begin{equation*}
= (-2)^{n-2} e^{-2x} \left[ 4(x^2 - x) - 2n(2x - 1) + n(n-1) \right]
\end{equation*}


            	
\begin{equation*}
= (-2)^{n-2} e^{-2x} \left[ 4x^2 - 4(n+1)x + n^2 + n \right]
\end{equation*}


代入$x=0, n=100$,    
\begin{equation*}
f^{(100)}(0) = (-2)^{98}\cdot [0-0+100^2+100]=2^{98}\cdot 10100
\end{equation*}

\end{solbox}

\textcolor{DeepPink}{方法D:[通项公式变形]}

\begin{solbox}[分步推导与答题规范]
理论公式:

$(e^{ax})^{(n)} = a^n e^{ax}$	$\left(\frac{1}{ax+b}\right)^{(n)} = \frac{(-1)^n n! a^n}{(ax+b)^{n+1}}$	$[\ln(ax+b)]^{(n)} = \frac{(-1)^{n-1} (n-1)! a^n}{(ax+b)^n} (n \ge 1)$

$(\sin ax)^{(n)} = a^n \sin\left(ax + \frac{n\pi}{2}\right)$		  $(\cos ax)^{(n)} = a^n \cos\left(ax + \frac{n\pi}{2}\right)$

操作(A):有理分式裂项:$\frac{1}{x^2-1}=\frac{1}{2}\left(\frac{1}{x-1} - \frac{1}{x+1}\right)$

操作(B):三角和差化积:$\sin^2 x$或$\sin x \cos 2x$,先降幂/化和差
\end{solbox}

\begin{thmbox}[题型二:差商与连续\&可导判定[注意分段讨论]]

可导:对于若已知$f(x)$在$x_0$处可导，则对称差商极限必定存在且等于$f'(x_0)$;\textcolor{red}{否则必须需要一个定点}

对称差商:比如已知极限$L = \lim\limits_{h \to 0} \frac{f(x_0 + h) - f(x_0 - h)}{2h}$存在

判断连续:$\lim\limits_{x \to x_0^+} f(x)=\lim\limits_{x \to x_0^-} f(x)= f(x_0)$				   (极限值=原函数函数值)

判断可导:
\begin{equation*}
\lim_{x \to x_0^+} \frac{f(x) - f(x_0)}{x - x_0}= \lim_{x \to x_0^-} \frac{f(x) - f(x_0)}{x - x_0}=A(\text{而非}\infty)
\end{equation*}
		 (左差商=右差商)

连续可导:$\lim_{x \to x_0^+} \frac{f(x) - f(x_0)}{x - x_0}= \lim_{x \to x_0^-} \frac{f(x) - f(x_0)}{x - x_0} = f'(x_0)$		    (先可导,然后判断导函数连续)

对称值:结论[1]:连续与绝对值连续

	  $f(x)$在$x_0$处连续$\implies |f(x)|$在$x_0$处必连续

	  $|f(x)|$在$x_0$处连续$\centernot\implies f(x)$在$x_0$处连续

      结论[2]:可导与绝对值可导

	  $f(x)$在$x_0$处可导$\centernot\implies |f(x)|$在$x_0$处可导

	  $|f(x)|$在$x_0$处可导$\centernot\implies f(x)$在$x_0$处可导

      结论[3]:函数值判断可导(设$f(x)$在$x_0$处可导)

	   [$f(x_0) \ne 0$]:        $|f(x)|$必可导

	   [$f(x_0) = 0$且$f'(x_0) = 0$]:$|f(x)|$必可导，且导数值为0

	   [$f(x_0) = 0$且$f'(x_0) \ne 0$]:$|f(x)|$必不可导(穿x轴的尖点)

      结论[4]:不可导的4种典型情形

	   断点:函数在$x_0$处不连续(不连续必不可导) 

	   尖点:左右导数均存在但不相等($f'_+(x_0) \ne f'_-(x_0)$，如$y = |x|$在$x=0$) 

	   垂直切线:差商极限为无穷大($\lim\limits_{x \to x_0} \frac{\Delta y}{\Delta x} = \infty$,如$y = x^{\frac{1}{3}}$在$x=0$)

	   高频震荡:差商极限震荡不存在(如$y = x \sin\frac{1}{x}$在$x=0$)

Eg:
\begin{equation*}
\vert{}f(x)\vert{} \le 1 - \cos x
\end{equation*}
,则f(x)[没有显式式子]在x=0处:

x=0带入:
\begin{equation*}
\vert{}f(0)\vert{} \le 1 - \cos 0 = 1 - 1 = 0
\end{equation*}
,所以$|f(0)|=0$

又
\begin{equation*}
-(1 - \cos x) \le f(x) \le 1 - \cos x
\end{equation*}
:

  左侧:$\lim\limits_{x \to 0} [-(1 - \cos x)] = -(1 - 1) = 0$

  右侧:$\lim\limits_{x \to 0} (1 - \cos x) = 1 - 1 = 0$

所以
\begin{equation*}
\lim_{x \to 0} f(x) =f(0)=0
\end{equation*}
,函数连续得证

又
\begin{equation*}
\big|\frac{f(x) - 0}{x - 0}\big| \leq \big|\frac{1-cosx}{x}\big|
\end{equation*}
,且$-\lim\limits_{x \to 0^+} \frac{\frac{1}{2}x^2}{x} = 0$,$\lim\limits_{x \to 0^-} \frac{\frac{1}{2}x^2}{x} = 0$,($f'_+(0) = f'_-(0) = 0$,所以$f'(0) = 0$)

\end{thmbox}
\begin{thmbox}[题型三:导数的几何意义]

Eg:设曲线$y = x^n$在点$(1, 1)$处的切线与$x$轴交点为$(\xi_n, 0)$,求$\lim\limits_{n \to \infty} f(\xi_n)$。

  求导得切线斜率:$y' = n x^{n-1} \implies k = y'(1) = n$

  写出切线方程:$y - 1 = n(x - 1)$

  求$x$轴交点:令$y = 0 \implies 0 - 1 = n(x - 1) \implies \xi_n = 1 - \frac{1}{n}$

  计算极限目标:
\begin{equation*}
\lim_{n \to \infty} f(\xi_n) = \lim_{n \to \infty} (\xi_n)^n = \lim_{n \to \infty} \left(1 - \frac{1}{n}\right)^n = e^{-1} = \frac{1}{e}
\end{equation*}


\end{thmbox}

\begin{equation*}
\Delta f(1) - df(1) = o(\Delta x)
\end{equation*}

\begin{equation*}
\lim_{\Delta x \to 0} \frac{\Delta f(1) - df(1)}{\Delta x} = \lim_{\Delta x \to 0} \frac{o(\Delta x)}{\Delta x} = 0
\end{equation*}

\begin{equation*}
\lim_{\Delta x \to 0} \frac{\Delta f(1) - df(1)}{\Delta x} = \lim_{\Delta x \to 0} \frac{f(1+\Delta x) - f(1)}{\Delta x} - \lim_{\Delta x \to 0} \frac{f'(1)\Delta x}{\Delta x} = f'(1) - f'(1) = 0
\end{equation*}
\begin{thmbox}[题型四:微分和导数]


\begin{equation*}
\Delta y = {A \Delta x}\text{（线性主部） }+ {o(\Delta x)}\text{高阶误差}
\end{equation*}
($dy = {A dx}\text{（线性主部） }+ {o(dx)}\text{高阶误差})$

  $\Delta y$(真实增量):函数在曲线上真实的纵坐标变化量，$\Delta y = f(x_0 + \Delta x) - f(x_0)$

  $dy$(微分/线性主部):切线上的纵坐标变化量，即 $dy = A \Delta x$

  在一元函数中，系数$A$必定就是导数，即 \textit{}$A = f'(x_0)$\textit{},
\begin{equation*}
dy = f'(x_0) \Delta x = f'(x_0) dx
\end{equation*}


\textbf{题型[1]微分逆推导数}

设函数$f(x)$在任意$x$处的增量为 $\Delta y = \frac{y \Delta x}{x + \sqrt{x^2 + y^2}} + o(\Delta x)$，且$f(0) = 1$，求 $y = f(x)$ 在 $x = 0$ 处的微分 $dy$。

  线性主部系数即导函数：$f'(x) = \frac{y}{x + \sqrt{x^2 + y^2}}$

  代入 $x = 0, y = f(0) = 1$：
\begin{equation*}
f'(0) = \frac{1}{0 + \sqrt{0^2 + 1^2}} = 1
\end{equation*}


则在$x = 0$处的微分为：
\begin{equation*}
dy = f'(0) dx = 1 \cdot dx = dx
\end{equation*}


\textbf{题型[2]复合函数的微分与线性主部}

设$y = f(x^2)$，当自变量$x$在$x = -1$处取得增量$\Delta x = -0.1$时，相应函数增量$\Delta y$的线性主部为$0.1$,求$f'(1)$。

  对复合函数求微分：
\begin{equation*}
dy = d[f(x^2)] = f'(x^2) \cdot d(x^2) = 2x f'(x^2) dx = 2x f'(x^2) \Delta x
\end{equation*}


  将$x = -1, \Delta x = -0.1, dy = 0.1$ 	代入：
\begin{equation*}
0.1 = 2(-1) \cdot f'((-1)^2) \cdot (-0.1)
\end{equation*}


  
\begin{equation*}
0.1 = 0.2 \cdot f'(1) \implies f'(1) = \frac{0.1}{0.2} = 0.5
\end{equation*}


\textbf{题型[3]:增量与微分的差商极限}

设$f(x)$在$x = 1$处可导,$\Delta f(1)$是$f(x)$在增量$\Delta x$下的函数增量，求极限$\lim\limits_{\Delta x \to 0} \frac{\Delta f(1) - df(1)}{\Delta x}$

\textbf{[1]} 根据高阶无穷小定义

由于 $f(x)$ 在 $x=1$ 处可导即说明可微，由微分定义：


因此：


\textbf{[2]} 拆项极限法


\end{thmbox}

\section{备用}

\textbf{1.泰勒公式}


\begin{equation*}
f(x) = \sum_{k=0}^{\infty} \frac{f^{(k)}(0)}{k!} x^k = f(0) + f'(0)x + \cdots + \frac{f^{(n)}(0)}{n!} x^n + \cdots
\end{equation*}

\textbf{2.莱布尼茨公式}


\begin{equation*}
(uv)^{(n)} = \sum_{k=0}^{n} C_n^k u^{(n-k)} v^{(k)} = u^{(n)}v + n u^{(n-1)}v' + \frac{n(n-1)}{2!} u^{(n-2)}v'' + \cdots + u v^{(n)}
\end{equation*}


\end{document}
